% ===========================================================================
% SOURCES: docs/claude_project/08_OPEN_ITEMS.md §2 (what was built, and when)
%          docs/design_decisions.md: DD-030, DD-031, DD-032, DD-033
%          hardware/rtl/axi/wgt_dma.v, axi_lite_csr.v
%          hardware/rtl/top/accel_sys.v, accel_axi.v
%          scripts/make_bd.tcl, scripts/run_bd.sh, scripts/make_sdcard.py
%          firmware/** (driver, preprocessing, tests, stubs)
%          build/bd_187p5_explore/
%          docs/lab_notebook.md: 2026-10-07 (the first board run),
%            2026-10-20 (the camera path)
%          docs/design_decisions.md: DD-035, firmware/linux/README.md
%
% This was §6.8 - one section - in the first outline. It is now a chapter
% because the work is: an AXI4 master, a CSR block, a board design that closes
% with a bitstream, a bare-metal application, and a preprocessing path that is
% bit-exact on the host. That is the Core success tier (DD-004) - REACHED on
% the board on 2026-10-07 (Section~\ref{sec:board-remaining}). The live camera
% (\DD{035}) is the Stretch tier - REACHED on 2026-10-22 - and closes the
% chapter (Section~\ref{sec:camera-path}).
% ===========================================================================
\chapter{System Integration and the Embedded Application}
\label{chap:integration}

\section{From an Out-of-Context Core to a System}
\label{sec:ooc-to-system}
% Open with the problem, because it explains every design choice after it:
% accel_top has 846 port bits against the package's ~300 PL user I/O, which is
% WHY every early number was out of context. Most of those bits were the load
% side, and they were pins only because nothing drove them.
% The progression, as a table:
%     accel_top  846 port bits
%     accel_sys  493   (wgt_dma absorbed 672 of them)
%     accel_axi  313   (196 AXI4 master + 114 AXI-Lite + 3 clock/reset/irq)
% accel_axi has no non-AXI ports at all.

\section{Weight DMA and AXI Interfaces}
\label{sec:dma}
% \DD{030}. Read-only AXI4 master, 128-bit RDATA, INCR bursts split at the
% 4 KB boundary and capped at 256 beats, four outstanding ARs.
% Two results:
%   - proven by accel_dma_tb: the whole \thInstructions-op network, \thElements
%     elements bit-exact, with NO testbench task between ld_req and the weight
%     buffers;
%   - it is FASTER than the task it replaces - \thDmaLoadCycles{} load cycles
%     against \thTaskLoadCycles{} - because it sustains one bank write per cycle.
% T-7 closed with it: ld_pchan carried three meanings, and ld_op is the fix.

\section{AXI-Lite Control Registers}
\label{sec:csr}
% \DD{031}. 13 registers, a 512-word program aperture, an irq line, six SLVERR
% guards. Proven by accel_axi_tb driving the entire network over nothing but
% two AXI interfaces, result read back over AXI-Lite.
% Cost: \thLutCsr{} LUT, \thFfCsr{} FF, no DSP and no BRAM.

\section{The Image Path}
\label{sec:image-path}
% \DD{032}. Weights, parameters AND the image on one read master - they are
% never concurrent, so there is no arbiter. HWC in DDR, 1.00 pixels/cycle,
% about 0.25 ms.
% Worth a sentence: axi_lite_csr_tb REGRESSED on the image ports and was right
% to - a good miniature of Section~\ref{sec:bit-exact-contract} in action.

\section{Constraints Without Invented Numbers}
\label{sec:xdc}
% A short section, and a real methodological point. accel_top_ooc.xdc asked to
% "tighten these once the DMA is real"; in accel_sys_ooc.xdc the load ports are
% not tightened, they are GONE - internal register-to-register paths with no
% input delay and no exception. That is strictly stronger than the 0.5 ns
% budget they had, because a budget is an assumption about a launch clock and a
% register-to-register path is a measurement.
% After \DD{032}, accel_axi_ooc.xdc is the first constraint file in the project
% with no invented number in it.

\section{Board-Level Implementation}
\label{sec:board-impl}
% scripts/make_bd.tcl: ZCU104 block design - PS + DDR4 + SmartConnect +
% accel_axi. Closes at \MHz{\thClkBoard} with WNS +\thWnsBoard{} ns, and
% produces a bitstream and an XSA.
% \thLutBoard{} LUT (\thLutPctBoard\%), \thBramBoard{} BRAM, \thUramBoard{} URAM,
% \thDspBoard{} DSP.
%
% Two findings no out-of-context run could have produced, and they are the
% reason this section exists:
%   1. THE PL CLOCK GRID IS COARSE. 1500 MHz over an integer divider gives only
%      150 / 166.7 / 187.5 / 214.3 / 250 MHz, and requesting 200 yields 187.5
%      SILENTLY. Quote \MHz{\thClkBoard}, the constrained clock, not the
%      achievable \MHz{\thFmaxBoard} - see Section~\ref{sec:cad-flow}.
%   2. THE CRITICAL PATH MOVED. In context the top three paths are wgt_dma's
%      output register into the weight buffers at zero logic levels - pure
%      route - not act_buffer's URAM (Section~\ref{sec:dd023}).
% Against accel_axi alone post-route (\MHz{\thFmaxAccelAxi}), the full system
% costs about 8% of the clock: much less than "the board is the unknown" implied.

\section{The Processing System Application}
\label{sec:ps-app}
% firmware/. The driver (accel_drv.c), the CSR map (accel_csr.h), the image
% source (img_src.c), the platform layer (plat_zynq.c).
% The sequence: SD card -> DDR -> PL -> result.
% Report that it is verified WITHOUT A BOARD: test_csr_map.c, test_drv_seq.c
% and test_preproc.c run on the host against stubs, and `make lint`
% compile-checks the bare-metal sources. Say what that does and does not buy
% (Section~\ref{sec:verification-limits}).
% Then that it has since run on one (Section~\ref{sec:board-remaining}), and
% that the two halves belong together: the host tests are why the first board
% run had nothing left to find in the driver.

\section{Preprocessing on the Target}
\label{sec:preproc-target}
% \DD{033}, from the hardware side; Section~\ref{sec:preproc} has the algorithm.
% The point to make here: this was the last unverified link in the chain. The
% \thElements-element bit-exactness began at 000_input.hex, so until this
% existed the thesis could not claim an image-to-class result even in
% simulation - only a tensor-to-class one.

\section{Staging the SD Card}
\label{sec:sdcard}
% scripts/make_sdcard.py stages the card with a KNOWN ANSWER, not just the
% blobs - so a board bring-up can fail loudly rather than plausibly. One
% paragraph; it is a small script with a good idea in it.

\section{Bring-Up on the ZCU104}
\label{sec:board-remaining}
\label{sec:bringup}
% WAS "What Remains for a Board Run". The board run happened, so the section
% keeps the PROCESS and the RESULT moves to
% Section~\ref{sec:bit-exactness-result}, as this note always said it would.
% (The old label is kept so the cross-references still resolve.)
%
% The run, from docs/lab_notebook.md 2026-10-07: the lab's ZCU104, booted
% over JTAG; image.bin on the SD card = the int8 healthy_hl6025 image staged
% by make_sdcard.py (Section~\ref{sec:sdcard}). The registers read back:
%     STATUS    0x3F22   done, result captured, program counter 63, no error bit
%     ARGMAX    0x6      class 2 (Healthy), valid
%     RESULT_LO 0xE7DFECAA, RESULT_HI 0xEA923B51
%               -> logits -4950 -6177 15185 -5486
% the same four integers as EXPECTED.txt. Bit-exact on silicon, first time.
%
% What it cost is worth one honest paragraph, because none of it was the
% design: every obstacle was in the Vitis 2022.2 flow on the development PC
% or on the bench - parallel BSP builds failing (make -j 1), xilffs with long
% file names in the standalone domain only (in the FSBL it overflows the OCM),
% a launch configuration that skipped the PS initialisation, a stale debugger
% connection, and an SD card not seated. The known answer is what made each
% of those quick to rule in or out: the accelerator was never a suspect.
%
% Then what is NOT yet recorded, stated as plainly: the RGB image through
% preproc.c (it must give the same four integers), the other three classes'
% cards, the UART image source, and booting from the SD card without JTAG.
% Under Linux on 2026-10-22 the same accel_drv.c and preproc.c did pass the
% RGB path and all four classes (Section~\ref{sec:camera-path}) - say so, and
% keep this list honest: for the bare-metal application itself they are
% still unrecorded. Fill them in from the bench; whatever is still missing at
% submission goes to Section~\ref{sec:future}.
% TODO(REDEPLOY): after DD-036's redeploy, repeat the known-answer runs on
% the new golden sets (one test-split photograph per class).

\section{The Live Camera Path under Linux}
\label{sec:camera-path}
% \DD{035} - the Stretch tier. A SECTION for now; if it grows past a few pages
% it can become a chapter of its own (left open in \DD{036}). RUN ON THE BOARD
% on 2026-10-22: all four known answers bit-exact under Linux, and the live
% camera classifying (docs/lab_notebook.md). The results belong to
% Section~\ref{sec:camera-result}; keep the design here, plus item 7 below.
%
% In this order, each the answer to a question a reader will ask:
%   1. WHY LINUX. A UVC webcam needs a USB host stack and a UVC driver; the
%      bare-metal BSP has neither. PetaLinux 2022.2 [xilinx_ug1144], built
%      from the ZCU104 BSP plus this design's XSA, provides both - and the
%      bitstream is still loaded by the FSBL from BOOT.BIN, so the FPGA design
%      does not change at all.
%   2. A FIFTH PLATFORM HOOK. The driver wrote (uint32_t)pointer into the base
%      registers - right only where a pointer is a physical address. Under
%      Linux it would point the DMA at an unrelated page and complete on
%      whatever was there: the failure shape this project designs out
%      everywhere else. accel_plat_phys translates all four blobs before any
%      register is written, and ACCEL_ERR_ADDR refuses what cannot be reached.
%   3. MEMORY THE DMA CAN READ. 16 MB of `no-map' reserved memory at
%      0x3800_0000, mapped through UIO - so no cache can hold a line of it,
%      stale data is impossible by construction, and the bare-metal cache
%      flush becomes a barrier. The cost, accepted: aligned 64-bit stores only,
%      and an uncached 150 KB copy per frame.
%   4. THE RESIZE. A camera frame is not 256x256, so Resize(256) is a real
%      resampling and DD-033's crop is not enough. camfront.c transcribes
%      Pillow's resampler [pillow], and is bit-exact against it on 15 frame
%      sizes. Of nine deliberate mistakes, eight are caught; the ninth changes
%      no byte, and that was settled by enumeration rather than argument. This is
%      Section~\ref{sec:preproc}'s method, extended to any frame.
%   5. WHAT STAYS OUTSIDE THE CHAIN, named: the camera's YUV-to-RGB conversion,
%      which has no reference to be exact against. A frame saved on the board
%      is saved after it, so check_camera_frame.py can re-derive the window,
%      the tensor and the logits from the RGB the board actually used.
%   6. THE SCREEN, in one paragraph: DisplayPort through fbdev, and the same
%      image over HTTP to a laptop for a projector. Neither is verified on the
%      board yet: there was no /dev/fb0 at boot on 2026-10-22.
%   7. WHAT THE BOARD RUN COST, one paragraph, as for the bare-metal bring-up
%      (Section~\ref{sec:board-remaining}) - and again none of it was the design:
%      the image was built on WSL2 (Section~\ref{sec:petalinux}); the kit's SD
%      card booted its own PetaLinux 2022.1 from an ext4 partition, so the
%      board ran the wrong system until it was booted without image.ub; the
%      UIO driver came out as a module despite the configuration fragment; the
%      card mounts root-only. The known answers made each quick to isolate:
%      the first run of the Linux port was bit-exact for all four classes.
% A figure belongs here (the camera-to-class pipeline across the PS and the
% PL); none exists yet - see thesis/README.md, Figures.
